\subsection{Exploratory extension: clean controls, a second monitor, and repeats}
\label{sec:extension}
After the original analysis, a separate locally frozen extension selected three additional upstream-test IDs per domain by a fixed hash, excluding all original held-out IDs. It used DeepSeek as actor and the same Sonnet customer, with native instructions or an added actor-only reminder to issue one call per message, and two unseeded runs per task/condition. This yields 24 new conversations on six new task IDs. The step limit was 90 in both extension conditions, versus 60 originally. No proposals were edited, serialized, rejected, or resampled to obtain compliance. A two-cell feasibility gate permitted expansion if either reminder trajectory was contract-clean; both feasibility cells remain in the allocation. This is exploratory evidence, not a new confirmatory test or a family-disjoint sample.

All 24 cells completed and replayed exactly. There are 20 state successes and 7 contract-clean state successes from 3 task IDs. Clean successes occur under both native and reminder instructions. Across the 12 task/condition repeat pairs, state success differs in 0 pairs and contract compliance differs in 3. These repetitions show why a single rollout should not be treated as a stable task property; they do not isolate a causal benefit of the reminder.

\begin{table}[t]
\centering\small
\caption{Separate extension cohort. Each row has three task IDs and two unseeded repeats. Clean success is state success without the measured interaction violation.}
\label{tab:extension}
\begin{tabular}{llrrr}
\toprule
Domain & Instructions & $n$ & State successes & Clean successes\\
\midrule
Retail & native & 6 & 6 & 2\\
Retail & reminder & 6 & 6 & 3\\
Airline & native & 6 & 4 & 1\\
Airline & reminder & 6 & 4 & 1\\
\bottomrule
\end{tabular}
\end{table}
DeepSeek V4.1 Flash also evaluated all 514 original public prefixes with the same broad prompt, schema, 512-token cap and 0.5 threshold, using its configured reasoning effort of none. There are 513 valid scores and 1 missing output(s). Of 510 matched valid Gemini/DeepSeek prefix pairs, 98 disagree at the threshold. Prefix pairs are descriptive comparisons, not independent samples. The fixed 0.5 threshold is uncalibrated for both models, so their score scales need not be comparable. In one matched batch proposal (retail task 70, DeepSeek renamed actor, event 6), DeepSeek explicitly describes the four-call violation but assigns risk 0.35, versus Gemini's 1.0. Alert disagreement therefore need not mean failure to recognize a constraint. An exploratory fixed-threshold grid is supplied without selecting a preferred threshold. Table~\ref{tab:expanded-timing} shows that the original coincidence between warning and violation onset does not generalize to the second monitor configuration. An alert preceding the interaction violation may concern another property; it is not thereby a correct forecast of eventual task failure.

Both monitors also score the new conversations. Gemini supplies 255 valid scores (1 missing), and DeepSeek supplies 256 (0 missing). Among the 7 clean state successes, the proxy-relative episode alert rate is 1/7 for Gemini and 5/7 for DeepSeek. These denominators now exist, unlike in the original cohort, but cover only 3 task IDs and do not certify semantic policy compliance. The native and reminder conditions remain separately tabulated in the artifact. Missing-score uncertainty is reported as lower/upper alert-count bounds rather than imputed negatives.

For Gemini, the descriptive whole-task bootstrap envelope is [0.0\%, 25.0\%], excluding 50 of 5,000 replicates with no clean-success denominator. For Deepseek, the descriptive whole-task bootstrap envelope is [0.0\%, 100.0\%], excluding 50 of 5,000 replicates with no clean-success denominator. These small-support intervals do not establish deployment calibration. DeepSeek shares a model family with half the original actors and all extension actors; this dependence is disclosed rather than treated as independent model validation.

The autonomous conversations are not comprehensively labeled for semantic-policy compliance. Optional human-review material is separate from the reported evidence; the later constructed eligibility experiment uses explicit simulator facts and executable labels instead. No independently human-reviewed result is claimed.
